\documentclass[11pt]{article}

% Packages
\usepackage[margin=1in]{geometry}
\usepackage{amsmath, amssymb}
\usepackage{graphicx}
\usepackage{hyperref}
\usepackage{listings}
\usepackage{xcolor}

% Code styling for Haskell / MLIR / Python
\lstset{
basicstyle=\ttfamily\small,
keywordstyle=\color{blue},
commentstyle=\color{gray},
stringstyle=\color{green!50!black},
breaklines=true,
frame=single
}

% Title Info
\title{Parsing Python with Haskell and Lowering to MLIR}
\author{Jegannarayanan Sivakasiulaganathan}
\date{\today}

\begin{document}

\maketitle

\begin{abstract}
	This report explores the design and implementation of a system that parses
	Python using Haskell and lowers it into MLIR. The motivation, architectural
	decisions, and results are discussed, along with use of AI and conclusions.
\end{abstract}

\tableofcontents
\newpage

% ----------------------------
\section{Introduction}
Brief overview of the project, goals, and motivation.

% ----------------------------
\section{MLIR: An Introduction}
\subsection{LLVM MLIR IR objects}
This section briefly looks at what MLIR IR objects look like, and why they are
important. In LLVM, MLIR is a intermediate representation between high-level
source code and LLVMIR:

\noindent Source Code (Python, C, C++, etc.) $\to$ MLIR $\to$ LLVMIR $\to$ Machine code

\noindent Some variants/dialects of MLIR:

\begin{enumerate}
	\item arith
	\item func
	\item linalg
\end{enumerate}

\noindent An example in arith:
\begin{verbatim}
module {
  func.func @add_example(%a: i32, %b: i32) -> i32 {
    %0 = arith.addi %a, %b : i32
    return %0 : i32
  }
}
\end{verbatim}

\noindent These representations allow for machine-agnostic standardized analyses and optimization passes.

\subsection{Core Concepts}
\begin{itemize}
	\item \textbf{Operations}: The basic unit of computation in MLIR. Each
	      operation represents a single step and can produce results and take
	      inputs.

	\item \textbf{Regions}: Containers attached to operations. Regions hold blocks
	      and allow for things like control flow constructs and
	      function bodies.
	\item \textbf{Blocks}: Ordered sequences of operations.
	      Blocks may take arguments, similar to function parameters.

	\item \textbf{Dialects}: Namespaces that group related operations and types.
	      They allow MLIR to support multiple domains Ex: \texttt{arith},
	      \texttt{func}, \texttt{llvm}.
\end{itemize}

\subsection{Why MLIR?}
\begin{itemize}
	\item \textbf{Multi-level abstraction}: MLIR allows different levels of
	      lowering, allowing us to reason on different levels of representations
	      (tensor, cf are examples)

	\item \textbf{Extensibility via dialects}: New domains can define custom
	      operations and types without modifying the core system.

	\item \textbf{Reusable infrastructure}: Optimization passes, analyses, and
	      transformations can be shared across different dialects and stages.
\end{itemize}
% ----------------------------
\section{Why Use Haskell to Parse Python}
\subsection{Motivation}
Why not Python or C++ to parse Python?

High school has algebraic types, and fluid composition. The type system and
monadic operations themselves lend to easier parsing. Furthermore, types like
`Either` and `Maybe` are very useful for error management in a fluid manner,
compared to how errors are managed and collected in imperative languages

\subsection{Parsing Strategy}
My strategy here was to leverage the library `language-python`, which parser
Python, and convert it into an AST Defined by the `language-python` itself.

Next, I wrote a converter to take the `language-python` AST and create my own
AST with all the operations I chose to support in Python, a much smaller subset
of Python. (for simplicity)

\subsection{Mapping to MLIR}
Lastly, my plan was to take this AST and lower it MLIR. I was planning to
leverage the library `mlir-hs` to do so.

I didn't fully succeed here.

What success would've looked like here is:

Taking my `Python--` AST and lowering this AST to an AST that mlir-hs
explicitly defined in its codebase.

Then my plan would have been to leverage the builder inside the project to
automatically emit MLIR code.

I stopped short at creating the MLIR AST converter.
% ----------------------------
\section{System Architecture}
\subsection{Pipeline Overview}
Python source $\rightarrow$ Haskell AST $\rightarrow$ MLIR AST $\rightarrow$
MLIR lowering.

\subsection{Example}
\begin{lstlisting}[language=Python]
x: int = 5
y: int = 10
z: int = x + y
\end{lstlisting}

\begin{lstlisting}

# Example MLIR (simplified)

%0 = arith.constant 5 : i32
%1 = arith.constant 10 : i32
%2 = arith.addi %0, %1 : i32
\end{lstlisting}

% ----------------------------
\section{AI Use}
\subsection{Use of AI}
I used AI heavily for the following:

Understanding Haskell language constructs as well as AST constructs and MLIR.
Generating code to my specifications Helping with LaTeX for writing

\subsection{Important statement about AI use}

While I used AI heavily to Generate code, I never generated code and looked the
other way. I always went through every line of code to understand both what the
intention of the code was at a higher level, as well as what the language
constructs actually were. I am confident in my ability to explain all of it.

Furthermore, I used AI To dig deep into both Haskell language constructs and
MLIR constructs.

% ----------------------------
\section{Results}
\subsection{What Works}
Currently, I am able to take an AST of a script defined like in
docs/planning/PythonMLIR--.md.

\subsection{What I learned}
Some of the things I learned in this project are:

* How to convert and walk ASTs using Haskell. In this regard, I learned a lot about how
* Deep dives on Haskell language constructs
* MLIR structure and operations: Operations, Regions, Blocks, Control Flow,
how if, for, function calls, etc all actually work under the hood. Terminators,
block args, and much more.
* How much more difficult a project is than it seems.

% ----------------------------
\section{Conclusion}
Summary and final thoughts.

% ----------------------------
\section{Future Work}
Possible improvements and extensions.

\end{document}

